\section{Erd\H{o}s Problem 71: cycle lengths in an arithmetic progression}

\subsection{1) Statement and approach}
Let $P=\{a+jb:j\geq0\}$, where $a,b$ are positive integers, and
suppose $P$ contains an even integer. Does a sufficiently large
constant average degree force a cycle with length in $P$?
Bollob\'as proved the answer is yes. We reconstruct a complete
deduction from a later, stronger theorem of Liu and Montgomery.
The latter theorem is an explicit external input, not reproved here.

\subsection{2) The external interval theorem}
Theorem 1.1 of Liu--Montgomery states that there is $d_0$ such that
every finite graph of average degree $d\geq d_0$ has some real number
\[
 \ell\geq\frac{d}{10(\log d)^{12}}
\]
for which every even integer in $[(\log\ell)^8,\ell]$ is a cycle
length of the graph. Logarithms are natural. Only the existence of
such an interval, and the fact that its endpoint tends to infinity
with $d$, will be needed.

\subsection{3) Fitting an even term of the progression}
The even terms of $P$ form another infinite progression
\[
 P\cap2\mathbb N=\{a_0+jh:j\geq0\},\qquad h=\operatorname{lcm}(b,2).
\]
Indeed, when $b$ is even, the hypothesis forces $a$ to be even and
$h=b$; when $b$ is odd, exactly every other term is even and $h=2b$.
Here $a_0$ denotes its first even term.

Choose $L$ so large that, for every $x\geq L$,
\[
 (\log x)^8\geq\max(a_0,3),\qquad
 x-(\log x)^8\geq h.
\]
Such an $L$ exists because $(\log x)^8=o(x)$ and grows without bound.
Choose $D\geq d_0$ so large that
$d/[10(\log d)^{12}]\geq L$ for all $d\geq D$.
The limit of this expression is infinite, so the choice is uniform
over all such $d$.

For a graph of average degree $d\geq D$, obtain $\ell\geq L$ from
the external theorem. The least member of $a_0+h\mathbb Z$ that is
at least $(\log\ell)^8$ is smaller than $(\log\ell)^8+h$ and at
least $a_0$. Hence it belongs to $P\cap2\mathbb N$ and is at most
$\ell$. The theorem supplies a cycle of exactly that length.
This proves the required assertion with a constant depending only
on the progression.

\subsection{4) Why the parity condition matters}
If $P$ contains only odd integers, complete bipartite graphs
$K_{r,r}$ have arbitrarily large average degree $r$ but no cycle
whose length is in $P$. Thus the parity assumption is necessary,
not an artifact of the interval theorem. The proof also handles
progressions beginning after any prescribed finite cutoff, since
the choice of $L$ incorporates $a_0$.

\textbf{Outcome: SOLVED, known affirmative resolution reconstructed.}
All deductions from the stated interval theorem are proved above;
its expansion and cycle-construction proof is the external dependency.
No novelty or independent formal verification is claimed.

\subsection{5) Sources}
\begin{itemize}
\item H. Liu and R. Montgomery, \emph{A solution to Erd\H{o}s and
Hajnal's odd cycle problem}, Theorem 1.1,
\texttt{https://arxiv.org/abs/2010.15802}.
\item T. F. Bloom, Problem 71,
\texttt{https://www.erdosproblems.com/71}, checked 23 September 2026.
\end{itemize}
The estimate refers to full coverage using the explicit external theorem.
\subsection{6) COMPLETION ESTIMATE}
\textbf{COMPLETION: 100\%}
